% Minimal retained preamble for S3 SIGMA 2022 article cards taken from
% sources/sigma-2022-012-a-quantum-0-infty-law/pinned-original-source.tex
% (doi:10.3842/SIGMA.2022.012, SIGMA 18 (2022), CC BY 4.0).
%
% Omitted and why: the publisher documentclass{sigma} block (source
% lines 1-7), the tikz packages (source lines 10-12, used only by the
% section-4 picture, which is not delivered) and the frontmatter
% commands - the sigma class is unavailable in the toolchain and the
% assembler supplies the article class. stmaryrd is kept: the author's
% intt macro (source line 25) uses \llbracket. allowdisplaybreaks is
% the author's own (source line 44).
% The author macro block (source lines 14-32, including the
% DeclareMathOperator for span) and the author theorem declarations
% (source lines 34-41, including the stray group brace on line 41) are
% copied verbatim below.
\usepackage{amsmath,amssymb,amsthm}
\usepackage{stmaryrd}
\allowdisplaybreaks
\DeclareMathOperator{\psan}{span}

\newcommand{\hilH}{\mathcal H}
\newcommand{\norm}[1]{\lVert #1\rVert}
\newcommand{\aver}[1]{\langle #1 \rangle}
\newcommand{\nN}{n\in \mathbb{N}}
\newcommand{\mN}{m\in \mathbb{N}}
\newcommand{\lN}{l\in \mathbb{N}}
\newcommand{\kN}{k\in \mathbb{N}}
\newcommand{\jN}{j\in \mathbb{N}}
\newcommand{\iN}{i\in \mathbb{N}}
\newcommand{\intt}[1]{\llbracket #1 \rrbracket}
\newcommand{\perpF}{{\perp _F}}
\newcommand{\expec}[1]{\mathbb{E}\left(#1\right)}
\newcommand{\algE}{{\mathcal E}}
\newcommand{\salg}{$\sigma$-algebra}
\newcommand{\rv}{random variable}
\newcommand{\Id}{\mathrm{Id}}
\newcommand{\ind}[1]{\mathbf{1}_{#1}}

\newtheorem{Theorem}{Theorem}[section]
\newtheorem{Corollary}[Theorem]{Corollary}
\newtheorem{Lemma}[Theorem]{Lemma}
\newtheorem{Proposition}[Theorem]{Proposition}
 { \theoremstyle{definition}
\newtheorem{Definition}[Theorem]{Definition}
\newtheorem{Example}[Theorem]{Example}
\newtheorem{Remark}[Theorem]{Remark} }
